% JETSPIN restarting documentation

\section{Restarting JETSPIN\label{sec:Restarting-JETSPIN}}

A simulation can be restarted after normal termination or error termination,
if the last was due to inadequate time allocation. In order to restart
a JETSPIN job the file \texttt{input.dat} is again required together
with the file \texttt{restart.dat}, which is the exact copy of the
\texttt{save.dat} file created by the previous job. In particular,
the \texttt{save.dat} file was written by JETSPIN at regular time
interval during the previous simulation. The file \texttt{save.dat}
is \textit{not} human readable, and it should be renamed as \texttt{restart.dat},
and located in the same directory of input file. If the user attempts
to restart JETSPIN without the \texttt{restart.dat} file, the job
will immediately exit with error. Note that the directive ``\textit{restart
yes}'' needs to be added in input file in order to restart a previous
JETSPIN job (see Section \ref{sec:Input-file}). Further, whenever
a job is restarted, it is also possible to reset the accumulated statistical
data, collected in the previous run, by adding the directive ``\textit{restart
reset}'' in input file. It is worth stressing that JETSPIN will append
new data to the files \texttt{statdat.dat},\texttt{ traj.xyz} , while
all the other output files will be \textbf{overwritten} (included
the \texttt{save.dat} file). Thus, the user is advised to keep backup
copies of the \texttt{restart.dat} file, noting the times it was written,
to help you avoid going right back to the start of a simulation.

The same restarting procedure can be also used for extending a simulation
beyond its initial allocation of timesteps. In this case, the directive
\textit{timesteps} should be adjusted in input file to the new total
number of steps required for the simulation. For instance, if a 10000
step simulation needs to be extended by a further 5000 steps, the
user should use the directive ``\textit{timesteps 15000}'' in the
input file and include the ``\textit{restart yes}'' directive. Further,
the restarting procedure can be useful in the event of machine failure.
Here, the job can be restarted in the same way from the \texttt{save.dat}
file, which is written at regular time intervals in order to meet
just such an emergency. We stress that the procedure is not foolproof
in the last event, since the job could be crash while the \texttt{save.dat}
file is being written.
